\documentclass[10pt,a4paper]{amsart}
\usepackage[margin=19mm]{geometry}
\usepackage{amsmath,amssymb,mathtools}
\usepackage[scr=rsfs,cal=euler]{mathalfa}
\usepackage{xfrac}
\usepackage[unicode,hidelinks,bookmarks=false]{hyperref}
\newcommand{\ie}{i.e.\ }
\newcommand{\cf}{cf.\ }
\newcommand{\resp}{respectively\ }
\newcommand{\Subsection}{\subsection}
\providecommand{\nobreakdash}{\nobreak\hbox{-}}
\setlength{\emergencystretch}{2em}
\multlinegap0pt
\usepackage{bm}
\usepackage{bbm}
\usepackage{dsfont}
\usepackage{xspace}
\usepackage{cancel}
\usepackage{mathtools}
\usepackage{amsthm}
\makeatletter
\@ifundefined{lemma}{\newtheorem{lemma}{Lemma}}{}
\@ifundefined{proposition}{\newtheorem{proposition}{Proposition}}{}
\@ifundefined{theorem}{\newtheorem{theorem}{Theorem}}{}
\makeatother
\newcommand{\Acc} {\operatorname{Acc}}
\newcommand{\nat} {\mathbb{N}} 
\title[On Expansions of Monadic Second-Order Logic with Dynamical P]{On Expansions of Monadic Second-Order Logic with Dynamical Predicates}
\author{Joris Nieuwveld, Jo\"el Ouaknine}
\date{Source: arXiv:2507.16581v1 (CC BY 4.0)}
\begin{document}
\maketitle

\noindent\emph{Source and licence.} On Expansions of Monadic Second-Order Logic with Dynamical Predicates. Version arXiv:2507.16581v1 (CC BY 4.0), distributed under the Creative Commons Attribution 4.0 International licence, as recorded in the arXiv licence metadata for this exact version and shown on the abstract page https://arxiv.org/abs/2507.16581v1. Full rights notice: licenses/arxiv-2507.16581-CC-BY-4.0.txt.\par\smallskip

\noindent\textit{Editorial scope.} Counted unit: the Proposition whose source label is \texttt{prop : reduction one unary predicate} in the article (source0.tex lines 557--561, source label \texttt{prop : reduction one unary predicate}), delivered together with the article's own complete original proof of it (source0.tex lines 562--616, one \texttt{proof} environment of 55 lines). No mathematics is written, repaired or re-derived: statement and proof are the article's own text line for line. Required context retained uncounted: thm : transducer reduction (lines 517--519); eq : characteristic word P (lines 530--532); thm : Buchi MSO char word (lines 548--550). Every definition, display and label of those lines is retained unchanged. Omitted from this package and not redistributed: 1--516, 520--529, 533--547, 551--556, 617--1654. Nothing inside a retained excerpt is omitted or altered: every retained body is delivered exactly as the article's lines are written, so no omission receipt of any kind accompanies this package. Citations: the retained excerpts carry no citation command at all, so no bibliography entry of the article is transcribed and no reference is redistributed. Reference adaptation: none was needed and none was made; every reference inside the delivered unit resolves inside it, so no unresolved reference or citation is produced. Printed slips kept literal: no printed word, symbol, hypothesis or number of the counted statement or of its proof was altered. Layout and class: the article's own class is \texttt{lipics-v2021} with packages including cite, amsmath, amssymb, amsfonts, algorithmic, graphicx, textcomp, xcolor, amsthm, tikz, hyperref, todonotes; the wrapper is the lane's pinned \texttt{amsart} editorial wrapper at 10pt on A4, which restates the theorem-like environments the retained text needs and adds the article's own macro definitions that the retained text needs (\texttt{\textbackslash Acc}, \texttt{\textbackslash nat}), with extra wrapper packages (bm, bbm, dsfont, xspace, cancel, mathtools). The delivered numbering is the unit's own. Assets: no retained span contains a figure environment, a graphics-inclusion command, a PDF-page inclusion, a table or a TikZ picture; no image file is copied into this package, the package declares no asset dependency and no image file is redistributed with it. Licence: version arXiv:2507.16581v1 of this article is distributed under the Creative Commons Attribution 4.0 International licence, as recorded in the arXiv licence metadata for this exact version and shown on the abstract page https://arxiv.org/abs/2507.16581v1. A full rights notice is delivered at licenses/arxiv-2507.16581-CC-BY-4.0.txt. Characters: every retained excerpt is pure ASCII, so the delivered engine, XeLaTeX with its default Latin Modern text font, reports no missing character.
\begin{lemma}\label{thm : transducer reduction}
Let $w \in \Sigma^\omega$ and $\mathcal{B}$ be a deterministic finite transducer. Then the problem $\Acc_{\mathcal{B}(w)}$ reduces to $\Acc_{w}$.
\end{lemma} 

\begin{equation}\label{eq : characteristic word P}
    w = 0^{p_0}10^{p_1 - p_0-1}10^{p_2 - p_1-1}10^{p_3 - p_2-1} \cdots \, .
\end{equation}

\begin{theorem}\label{thm : Buchi MSO char word}
    The decidability of the MSO theory of the structure $\langle \nat ; < , P \rangle$ is Turing equivalent to $\Acc_w$, where $w$ is the characteristic word of $P$.
\end{theorem}

\begin{proposition}\label{prop : reduction one unary predicate} 
Let $P \subseteq \nat$ be an infinite, recursive, and effectively
sparse predicate with enumeration $\langle p_m\rangle_{m=0}^\infty$.
If, for each $M \ge 2$, $\Acc_{\langle p_m\bmod M\rangle_{m=0}^\infty}$ is decidable, the MSO theory of the structure $\langle \nat; <, P \rangle$ is decidable.
\end{proposition}

\begin{proof}
  By Thm.~\ref{thm : Buchi MSO char word}, it is sufficient to be able
  to decide whether a deterministic Muller automaton
  $\mathcal{A} = (\{0,1\}, Q, q_{\text{init}}, \delta, \mathcal{F})$
  accepts the characteristic word \eqref{eq : characteristic word P}.
  Let us restrict $\mathcal{A}$ to a directed graph $G$ with nodes $Q$
  and $0$-transitions as arrows.

  By construction, every node in $G$ has outdegree $1$ and so each
  state is in at most one cycle in $G$.  We can therefore compute the
  least common multiple of the cycle lengths in $G$ (call this number
  $M$) and the longest path to a cycle (call this number $N$).  Then,
  for all states $q$ and numbers $n \ge M + N$ and $d \ge 1$, reading
  $0^n$ and $0^{n+dM}$ leads to journeying through the exact same set
  of states and ending up in the same state.

  Let $K$ be such that $p_{m+1} - p_m \ge M + N+1$ for all $m \ge K$
  (which can be computed as $P$ is effectively sparse).  We construct
  a deterministic finite transducer $\mathcal{B}$ that hard-codes the
  initial segment
  $w_\text{init} := 0^{p_0}10^{p_1-p_0-1}1\cdots
  0^{p_{K+1}-p_{K}-1}1$.  For $m \ge K$, after reading
  $(p_m\bmod M)$ and $(p_{m+1}\bmod M)$, $\mathcal{B}$ outputs
  $0^{k_m}1$, where $M+N < k_m \le 2M + N$ is congruent to
  $p_{m+1} - p_m - 1$ modulo $M$.  Then, by construction, a state $q$
  is visited infinitely often upon reading the characteristic word of
  $P$ if and only if $q$ is visited infinitely often when
  $\mathcal{A}$ reads
  $w_\text{init}\mathcal{B}(\langle p_m\bmod M\rangle_{m=K+1}^\infty)$.

    Recall that $\Acc_{\langle p_m\bmod M\rangle_{m=0}^\infty}$ is assumed to be decidable. 
    Then $\Acc_{\langle p_m\bmod M)_{m=K+1}^\infty}$ is also
    decidable (by hard-coding the initial segment) and thus
    $\Acc_{\mathcal{B}(\langle p_m\bmod M\rangle_{m=K+1}^\infty)}$ is decidable by Thm.~\ref{thm : transducer reduction}.
    Therefore $\Acc_{w_\text{init}\mathcal{B}(\langle p_m\
      \mathrm{mod}\ M\rangle_{m=K+1}^\infty)}$ is decidable (by again
    hard-coding the initial segment), and by construction, $\mathcal{A}$ accepts
    the characteristic word of $P$ if and only if $\mathcal{A}$
    accepts $w_\text{init}\mathcal{B}(\langle p_m\bmod M\rangle_{m=K+1}^\infty)$. 
    Hence $\Acc_w$ is decidable, as required.
    %
    % Using Theorem~\ref{thm : contraction method}, we obtain the numbers $N$ and $M$.
    % As $P$ is effectively sparse, we can hard code a finite prefix, shift the predicate, and assume that $p_0 = 0$ and $p_{m+1} - p_m > MN$ for all $m \ge 0$.
    % Then, by Theorem~\ref{thm : contraction method}, $\mathcal{A}$ accepts $w$ if and only if $\mathcal{A}'$ (with a new initial state derived from the hard-coded prefix) accepts
    % \begin{equation*}
    %     w' = 10^{NM + (p_1-p_0-1 \pmod M)}10^{NM + (p_2-p_1-1 \pmod M)}\cdots.
    % \end{equation*}
    %
    % We define the automaton $\mathcal{A'} = (\{0,\dots,M-1\}, Q \times 2^Q, (q_\text{init}, \emptyset), \delta', \mathcal{F}')$ by setting $\delta'((q, A), m) = (q', A'))$ such that when the state $q$ in $\mathcal{A}$ reads $10^{NM+m}$, it visits the states $A'$ and ends up in the state $q'$.
    % Moreover, $\mathcal{F}'$ as all subsets of $Q \times 2^Q$ whose union of their second coordinates equals $\mathcal{F}$.
    %
    % A state $q$ in $\mathcal{A}$ is visited infinitely often if and only if there is a state $q'$ and a $0 \le m \le M-1$ such that infinitely often, when $\mathcal{A}$ is in $q'$, it will read the word $10^{NM+m}1$ of $w'$ next and passes by $q$ while reading this word.
    % Hence, $\mathcal{A}$ accepts $w'$ if and only if $\mathcal{A}'$ accepts \eqref{eq : modulo char word}.  
    % The result now follows, as we can decide whether $\mathcal{A}'$ accepts \eqref{eq : modulo char word} by the hypothesis.
\end{proof}

\end{document}
